\section{Solving the three-dimensional boundary deformation}
\label{n3:existence:section}

The boundary estimates have bounded $f$, $Z$ and $Df$ for every
admissible solution of the continuation equations. We now turn those
bounds into existence. Two points need attention. The divergence
equation for $Z$ contains $|D^2f|^2$, so the bounded variables do not
immediately give a classical elliptic estimate. Also, a compact
solution map must be defined on an open set of arbitrary inputs,
before imposing the admissibility inequality $t>-\epsilon$.

We address these points in order. A gradient-aligned scalar equation
gives both continuity of $Df$ and a local integral bound for $D^2f$.
An oscillation argument then gives continuity of $Z$, after which
ordinary boundary regularity applies. Separately, a scalar Dirichlet
problem can be solved for every bounded input $Z$, even when its
associated $t$ lies below the floor. Freezing that scalar solution in
the divergence equation produces the compact map. The floor exclusion
and the preceding estimates keep its fixed points inside the
admissible set along the homotopy of
Proposition~\ref{n3:deformation:homotopy}.

All regularity constants in this section may depend on the fixed
deformation parameters $N$, $\epsilon$ and $\tau$. They are not used
in any earlier argument requiring polynomial dependence on $N$.
The estimates needed for exhaustion will be uniform in the outer
radius $R$ at fixed $N$.

\subsection{From a bounded gradient to a Hessian integral bound}

The first equation has two equal transverse eigenvalues; its remaining
eigenvector is the gradient of its solution. This structure gives a
gradient estimate before the axial eigenvalue is continuous. The
accompanying Hessian estimate will control the source in the second
equation.

In the next two lemmas, $D$ denotes coordinate derivatives; it is not
the scalar $1+l^2|\nabla f|^2$ from the deformation.  Balls
are Euclidean coordinate balls, and a boundary ball is their
intersection with one side of a smooth boundary.  A family of boundary
patches has bounded geometry here if normal-coordinate charts, their
inverses, the metric and its inverse, and the derivatives through the
orders used below have fixed bounds.  We only apply the lemmas to
smooth metrics and smooth boundaries.  Their doubled metrics will be
Lipschitz and smooth on each side of a single plane.

\begin{lemma}[Axial gradient estimate]\label{n3:existence:gradient}
Let $U'$ be compactly contained in an interior or boundary coordinate
patch $U$ in dimension three.  In a boundary patch assume that $z$ is
constant on the boundary face.  Suppose $z$ is $C^3$, up to that face when present,
$|\nabla z|\le L$, and, almost everywhere,
\begin{equation}\label{n3:existence:axial-equation}
 A_\chi:\Hess z=G,\qquad
 A_\chi=I-(1-\chi)e\otimes e,\qquad
 e=\frac{\nabla z}{|\nabla z|},\qquad
 0<\chi_0\le\chi\le1,
\end{equation}
where $\|G\|_\infty\le Q$.  At zero gradients any measurable unit
axis for which the equation holds is allowed.  Then there are
$\alpha\in(0,1)$ and $C<\infty$, depending only on $\chi_0$, the
patch geometry and the separation of $U'$ from the other patch
boundaries, such that
\begin{equation}\label{n3:existence:gradient-estimate}
 \|\nabla z\|_{C^{0,\alpha}(U')}
 \le C(L+Q).
\end{equation}
For all sufficiently small coordinate balls centered in $U'$,
with intersection with the domain understood at a boundary, one also
has
\begin{equation}\label{n3:existence:hessian-morrey}
 \int_{B_r\cap U}|\Hess z|^2\dd V_g
 \le C(L+Q)^2r^{1+2\alpha}.
\end{equation}
In particular, the constants are independent of continuity moduli of
$\chi$ and $G$.
\end{lemma}

\begin{proof}
We first obtain an $L^p$ Hessian estimate for some $p>2$.
A divergence identity then shows that a gradient nearly attaining its
upper bound forces the solution to be nearly affine. The equation
with a fixed gradient direction improves that affine approximation.
The same steps survive reflection at a constant Dirichlet face,
and iteration gives the stated estimate.

\paragraph{A Cordes estimate above exponent two.}
For a Euclidean unit vector $e$,
\begin{equation}\label{n3:existence:cordes-defect}
 \|I-A_\chi\|_F=1-\chi\le1-\chi_0.
\end{equation}
For a compactly supported smooth function $v$, Plancherel's theorem
gives $\|D^2v\|_2=\|\Delta v\|_2$.  The Calder\'on--Zygmund
operator taking $\Delta v$ to the full array $D^2v$ is bounded on
$L^p$, and interpolation makes its norm arbitrarily close to one as
$p\to2$.  Choose once and for all
\begin{equation}\label{n3:existence:p-zero}
 2<p_0<3,\qquad
 C_{p_0}(1-\chi_0/2)<1.
\end{equation}
Rescale a metric-normalized patch until the coordinate principal
matrix differs from its Euclidean axial matrix by less than
$\chi_0/2$ in Frobenius norm.  Absorbing this difference in the
Laplacian estimate, applying a cutoff, and using
$\|Dv\|_p\le\eta\|D^2v\|_p+C_\eta\|v\|_p$ gives
\begin{equation}\label{n3:existence:local-cordes}
 \|D^2v\|_{L^{p_0}(B_{1/2})}
 +\|Dv\|_{L^{p_0}(B_{1/2})}
 \le C\bigl(\|v\|_{L^{p_0}(B_1)}
       +\|A_\chi:\Hess v\|_{L^{p_0}(B_1)}\bigr).
\end{equation}
The Christoffel terms are bounded first-order terms in this estimate.
The same argument applies when $v=z-L_0$ and $L_0$ is coordinate
affine: its extra right side is bounded by
$C\|Dg\|_\infty|DL_0|$.
The estimate uses bounded measurable principal coefficients only;
the assumed $C^3$ regularity allows all of its applications
without an approximation issue.  The linear estimates just used are
the ordinary Laplacian estimates; see
\cite[Chapters 7 and 9]{GilbargTrudinger2001}.

\paragraph{A nearly maximal gradient forces small Hessian energy.}
The preceding estimate does not yet give continuity of $Dz$. To gain
continuity we use the fact that the axial direction is $\nabla z$
itself. Put $H=\Hess z$ and $P=\nabla z$.  Equation~\eqref{n3:existence:axial-equation}
gives the pointwise vector identity
\begin{equation}\label{n3:existence:gradient-flux}
 A_\chi\nabla\frac{|P|^2}{2}-GP
       =(H-\Delta z\,g)P.
\end{equation}
Indeed $HP=|P|H(e,\cdot)$ and
$G=\Delta z-(1-\chi)H(e,e)$, so the two axial terms cancel.
The identity is also valid at $P=0$.  Taking the divergence of its
right side, rather than differentiating $\chi$, yields
\begin{equation}\label{n3:existence:bochner-flux}
 \Div\left(A_\chi\nabla\frac{|P|^2}{2}-GP\right)
 =|H|^2-(\Delta z)^2+\Ric(P,P).
\end{equation}
Choose $a_0>0$ so small that
$(1+a_0)(1-\chi_0)^2<1$.  Young's inequality and the equation show
\begin{equation}\label{n3:existence:positive-hessian}
 |H|^2-(\Delta z)^2
 \ge c_0|H|^2-C_0G^2,
\end{equation}
where $c_0>0$ depends only on $\chi_0$.

Consider solutions normalized by $|\nabla z|\le1$ on $B_1$.
After making the scale small, let the forcing and the rescaled
geometric errors have size at most $\delta$.  The nonnegative function
$s=1-|\nabla z|^2$ then satisfies
\begin{equation}\label{n3:existence:s-supersolution}
 \Div(A_\chi\nabla s+2G\nabla z)
 \le -2c_0|\Hess z|^2+C\delta
\end{equation}
in a weak sense.  In the Euclidean case the last term can be replaced
by $C\delta^2$.  We allow $C\delta$ to include the geometric errors.

For the compactness step, suppose
the normalized errors tend to zero in a sequence and
$\inf_{B_{1/2}}s\to0$.  The inhomogeneous weak Harnack inequality,
applied to Equation~\eqref{n3:existence:s-supersolution} after dropping
its favorable Hessian term, gives $s\to0$ in measure on $B_{1/2}$.
The boundedness $0\le s\le1$ also gives convergence in every finite
$L^p$ norm there.  To recover the Hessian information, take a cutoff
$\eta$ supported inside $B_{1/2}$ and use
$\eta^2(b-s)_+$ in the weak supersolution inequality.  Ellipticity
and Young's inequality give
\begin{equation}\label{n3:existence:low-truncation}
 \int_{\{s<b\}}\eta^2|Ds|^2
 \le Cb^2\int|D\eta|^2+o(1)
\end{equation}
for each fixed $b>0$ as the errors vanish.  The error fields in
Equation~\eqref{n3:existence:s-supersolution} are bounded and tend to
zero, so their products with $Ds$ are absorbed in deriving this
inequality.  Equation~\eqref{n3:existence:local-cordes} supplies a
uniform $L^2$ bound for $Ds$.  Its integral over $\{s\ge b\}$ tends
to zero in $L^1$ by Cauchy's inequality and convergence in measure.
On $\{s<b\}$ use Equation~\eqref{n3:existence:low-truncation}, then let
$b\downarrow0$.  Thus $Ds\to0$ locally in $L^1$.
Testing Equation~\eqref{n3:existence:bochner-flux} with a nonnegative
compact cutoff and using Equations~\eqref{n3:existence:positive-hessian}
and \eqref{n3:existence:gradient-flux} now gives
\begin{equation}\label{n3:existence:hessian-flat-limit}
 \int_{B_{1/4}}|\Hess z|^2\longrightarrow0.
\end{equation}
Uniform Lipschitz compactness implies that, after subtracting a
constant and taking a subsequence, $z$ tends uniformly on a smaller
ball to an affine function.  Its slope has length one because
$|\nabla z|\to1$ in measure and the Hessians tend to zero.

Consequently, for any prescribed sufficiently small $b_*>0$ there
are fixed $r_*\in(0,1/4)$, $k_*\in(r_*,1)$ and $\delta_*>0$ such
that every normalized solution with errors at most $\delta_*$ has
one of the following properties:
\begin{equation}\label{n3:existence:drop-or-flat}
 \begin{split}
 &\sup_{B_{r_*}}|\nabla z|\le k_*;\qquad\text{or}\\
 &\sup_{B_{r_*}}|z-L_0|\le b_*r_*,
       \qquad \tfrac12\le|DL_0|\le2.
 \end{split}
\end{equation}
If this conclusion failed, solutions whose gradient suprema approach
one and whose errors approach zero would contradict
Equation~\eqref{n3:existence:hessian-flat-limit}.  A gradient-drop
constant can always be weakened toward one to ensure $k_*>r_*$.

\paragraph{Regularity when the direction is fixed.}
The second alternative in \eqref{n3:existence:drop-or-flat} supplies
a nonzero affine slope. The comparison equation for that slope has
a fixed direction, although its axial coefficient remains merely
measurable. Let $e_0$ be a Euclidean unit vector and suppose
\begin{equation}\label{n3:existence:fixed-axis}
 \Delta_{e_0^\perp}Y_0+\chi(x)\partial_{e_0e_0}Y_0=0
 \quad\text{in }B_{3/4},\qquad
 \|Y_0\|_{W^{2,2}(B_{3/4})}\le C_1.
\end{equation}
The derivative $v=\partial_{e_0}Y_0$ belongs to $W^{1,2}$ and,
in distributions, satisfies
\begin{equation}\label{n3:existence:axial-derivative-divergence}
 \Delta_{e_0^\perp}v+
 \partial_{e_0}(\chi\,\partial_{e_0}v)=0.
\end{equation}
This follows by differentiating the distribution $\chi\partial_{e_0}v$
as a whole.  It does not assume a derivative of $\chi$ exists as a
function.  De Giorgi's estimate for this uniformly elliptic
divergence equation, followed by Caccioppoli applied to
$v-v(x_0)$, gives an exponent $\alpha_1\in(0,1)$ and
\begin{equation}\label{n3:existence:axial-derivative-morrey}
 [v]_{C^{0,\alpha_1}(B_{1/2})}\le C,
 \qquad
 \int_{B_r(x_0)}|Dv|^2\le Cr^{1+2\alpha_1}
\end{equation}
for balls in a smaller fixed interior region.  These are the scalar
divergence estimates with ellipticity $\chi_0$; see
\cite[Chapter 8]{GilbargTrudinger2001}.

All components of $DY_0$, not only $v$, are controlled by this
estimate.  In fact,
\begin{equation}\label{n3:existence:poisson-source}
 \Delta Y_0=f_0,\qquad
 f_0=(1-\chi)\partial_{e_0}v,
 \qquad
 \int_{B_r(x_0)}|f_0|\le Cr^{2+\alpha_1}.
\end{equation}
The last inequality follows from Cauchy's inequality and
Equation~\eqref{n3:existence:axial-derivative-morrey}.  Take the Newton
potential of $f_0$ restricted to a slightly larger interior ball;
its difference from $Y_0$ is harmonic in that ball.  For two points
at distance $h$, the contribution to the difference of gradient
potentials from annuli of radius $r\le2h$ is bounded by
$C\sum r^{\alpha_1}\le Ch^{\alpha_1}$.  On the remaining annuli
the difference of gradient kernels is at most $Ch r^{-3}$, so their
contribution is
\begin{equation}\label{n3:existence:potential-annuli}
 C h\sum_{r\ge2h}r^{\alpha_1-1}
 \le Ch^{\alpha_1}.
\end{equation}
The harmonic remainder has its usual interior derivative estimates.
Therefore, decreasing the interior ball if needed,
\begin{equation}\label{n3:existence:model-gradient}
 \|Y_0\|_{C^{1,\alpha_1}(B_{1/4})}\le C.
\end{equation}

\paragraph{Improving the affine approximation.}
We now compare the original equation with the fixed-direction equation.
Assume that, for a coordinate-affine function $L_0$,
\begin{equation}\label{n3:existence:flat-input}
 \|z-L_0\|_{L^\infty(B_1)}\le b,
 \qquad \tfrac14\le|DL_0|\le4,
 \qquad |\nabla z|\le8.
\end{equation}
Let the forcing and the first-derivative metric errors be at most
$b\delta$.  For $Y=(z-L_0)/b$, Equation~\eqref{n3:existence:local-cordes}
on fixed nested balls gives
\begin{equation}\label{n3:existence:normalized-residual}
 \|Y\|_{W^{2,p_0}(B_{7/8})}\le C.
\end{equation}
In particular, this estimate does not require a bound for
$DY$ in the supremum norm.  Since $Dz=DL_0+bDY$, the variable axis
converges in measure to $e_0=DL_0/|DL_0|$ as $b\to0$.  More
quantitatively, the coefficient error $E$ caused by replacing this
axis by $e_0$ is bounded and satisfies
\begin{equation}\label{n3:existence:axis-error}
 \|E\|_{L^q(B_{7/8})}\le Cb^{\vartheta}+Cb\delta,
 \qquad q=\frac{2p_0}{p_0-2},
 \qquad \vartheta=\frac{p_0-2}{2}>0.
\end{equation}
Indeed the unit-direction map is Lipschitz where
$|bDY|<|DL_0|/2$, and elsewhere its difference is bounded; the
$L^{p_0}$ bound in Equation~\eqref{n3:existence:normalized-residual}
and interpolation give Equation~\eqref{n3:existence:axis-error}.
The same estimate includes the metric principal-part error.
H\"older's inequality, with $1/2=1/q+1/p_0$, shows that
\begin{equation}\label{n3:existence:small-fixed-axis-source}
 \bigl\|[I-(1-\chi)e_0\otimes e_0]:D^2Y\bigr\|_{L^2(B_{3/4})}
 \le C(b^{\vartheta}+\delta).
\end{equation}

Solve for a zero-Dirichlet correction $w$ on $B_{3/4}$ with the
left side of Equation~\eqref{n3:existence:small-fixed-axis-source} as
its source.  This does not require a nondivergence Green function.
The convex Dirichlet Hessian inequality
$\|D^2w\|_2\le\|\Delta w\|_2$ and
Equation~\eqref{n3:existence:cordes-defect} make the Laplace solution
operator a contraction on $W^{2,2}\cap W^{1,2}_0$.  Hence
\begin{equation}\label{n3:existence:fixed-axis-correction}
 \|w\|_{W^{2,2}(B_{3/4})}
 \le C\chi_0^{-1}(b^{\vartheta}+\delta),\qquad
 \|w\|_\infty\le C(b^{\vartheta}+\delta).
\end{equation}
The last implication uses $W^{2,2}\hookrightarrow C^{0,1/2}$ in
dimension three.  The convex Hessian inequality is the
Miranda--Talenti estimate
\cite[Section~2, Theorem~3, p.~297]{Talenti1965}.
Now $Y_0=Y-w$ satisfies Equation~\eqref{n3:existence:fixed-axis} with
a uniform $W^{2,2}$ bound.

Choose $0<\alpha_2<\alpha_1$ and then a fixed
$\lambda_0\in(0,1/4)$ such that
$C\lambda_0^{1+\alpha_1}\le\lambda_0^{1+\alpha_2}/4$ in
Equation~\eqref{n3:existence:model-gradient}.  By first taking
$\delta$ small and then $b\le b_*$ small,
Equation~\eqref{n3:existence:fixed-axis-correction} has supremum norm
at most $\lambda_0^{1+\alpha_2}/4$.  Taylor expansion of $Y_0$
at the center therefore gives an affine function $L_1$ with
\begin{equation}\label{n3:existence:improved-flatness}
 \sup_{B_{\lambda_0}}|z-L_1|
 \le b\lambda_0^{1+\alpha_2},\qquad
 |DL_1-DL_0|\le Cb.
\end{equation}
All constants are uniform in the slope range in
Equation~\eqref{n3:existence:flat-input}.

\paragraph{Reflection at a constant Dirichlet face.}
Subtract the boundary value and use normal coordinates with the face
given by $x_3=0$.  Reflect the metric across this plane and reflect
$z$ oddly.  Tangential derivatives of $z$ vanish on the face, so
the extension is $C^1$.  There is no distributional Hessian atom.
The doubled metric is Lipschitz; the coefficient and lower-order
estimates in Equation~\eqref{n3:existence:local-cordes} hold with its
essential first-derivative bound.

The flux in Equation~\eqref{n3:existence:gradient-flux} may have a
normal jump.  On either side of the face its normal trace is
\begin{equation}\label{n3:existence:reflected-flux}
 \bigl((\Hess z-\Delta z\,g)\nabla z\bigr)_\nu
   =-(\tr_T\Hess z)\,\partial_\nu z.
\end{equation}
Because the boundary value is constant, $\tr_T\Hess z$ is its
mean-curvature coefficient times $\partial_\nu z$.  The jump is
therefore bounded by $C|\mathrm{II}|\,|\nabla z|^2$.  At scale
$r$ it is a bounded plane density of size $O(r)$.  Write the plane
density and flux after multiplication by the reflected coordinate
volume density, so that the following divergence is Euclidean.
Such a signed density $a(x')\,\delta_{\{x_3=0\}}$ is the divergence of the
bounded field $a(x')\mathbf 1_{\{x_3>0\}}\partial_3$.
Thus its negative part is another small divergence error in
Equation~\eqref{n3:existence:s-supersolution}.  Smooth-side curvature
errors are $O(r^2)$.  Weak Harnack and the low-truncation argument
remain valid, and the flatness comparison uses the same Lipschitz
metric estimates.  This proves all the preceding alternatives for
balls meeting the reflection plane as well.

\paragraph{Iteration and the required Hessian bound.}
There are now two ways to improve the solution on a smaller ball:
its gradient bound decreases, or its approximation by a nonconstant
affine function improves. We choose the constants once so that either
route can be repeated.
Keep $\alpha_2,\lambda_0$ from the affine improvement fixed.  Choose
$b_*$ sufficiently small that
$Cb_*/(1-\lambda_0^{\alpha_2})<1/4$, with $C$ the slope-increment
constant in Equation~\eqref{n3:existence:improved-flatness}.
Then choose $r_*,k_*,\delta_*$ in
Equation~\eqref{n3:existence:drop-or-flat} for this fixed $b_*$.
Normalize initially by a fixed multiple of $L+Q$ and choose a
uniformly small initial radius so that the metric and forcing errors
satisfy all the preceding smallness conditions.  If the first case
of Equation~\eqref{n3:existence:drop-or-flat} holds, rescale the ball
by $r_*$ and the gradient bound by $k_*$.  The normalized forcing
decreases by $r_*/k_*<1$.  Repeating this step either continues
forever or reaches the second case.  In the former case the gradient
decays with exponent $\log k_*/\log r_*$.
In the latter case apply Equation~\eqref{n3:existence:improved-flatness}
successively, retaining exponent $\alpha_2$.  At each affine step,
rescale $z$ by dividing by the spatial scale, subtracting only an
additive constant; this preserves its gradient bound.
Its slope changes have total size at most
$Cb_*/(1-\lambda_0^{\alpha_2})$, so the slopes remain in $[1/4,4]$.
At the $j$th step the error size is $b_*\lambda_0^{j\alpha_2}$,
while the normalized forcing and metric first-derivative errors have
size $O(\lambda_0^j)$.  Since $\alpha_2<1$, their ratio to the error
size only decreases.  This justifies every subsequent application
of the improvement lemma.  To combine the two branches, now choose
\begin{equation}\label{n3:existence:iteration-exponent}
 0<\alpha\le\min\left\{\alpha_2,
              \frac{\log k_*}{\log r_*}\right\},\qquad \alpha<1.
\end{equation}

At every center we have consequently obtained affine functions
$L_{x,r}$ satisfying
\begin{equation}\label{n3:existence:campanato-affine}
 \sup_{B_r(x)}|z-L_{x,r}|
 \le C(L+Q)r^{1+\alpha}.
\end{equation}
The slope differences at consecutive radii are bounded by
$C(L+Q)r^\alpha$; comparing the affine approximations on
overlapping balls gives Equation~\eqref{n3:existence:gradient-estimate}.
This also identifies their limiting slope with $Dz$, since $z$ is $C^1$.  The reflected version gives the estimate up to a boundary.

Finally apply the $L^2$ version of
Equation~\eqref{n3:existence:local-cordes} to $z-L_{x,2r}$ after
rescaling.  The affine error contributes
$C(L+Q)r^{1/2+\alpha}$ to the Hessian $L^2$ norm; the bounded
forcing and affine Christoffel term contribute $C(L+Q)r^{3/2}$.
Squaring and using $\alpha<1$ proves
Equation~\eqref{n3:existence:hessian-morrey}.
\end{proof}

\subsection{Continuity in the presence of a quadratic source}

The preceding lemma controls $|D^2f|^2$ by its mass on small balls.
We next show why that control suffices for the $Z$-equation.
An exponential change of variable absorbs its quadratic gradient
term. The remaining measure has a small Green potential, and weak
Harnack then decreases the oscillation. We allow a measure source
so that the argument also covers the inner boundary after reflection.

\begin{lemma}[Natural growth with a measure source]
\label{n3:existence:natural-growth}\label{b:lem:natural-growth}
Let $n\in\{3,4\}$ and let $Z\in L^\infty\cap W^{1,2}_{\mathrm{loc}}$
on a ball in $\R^n$. Fix a Borel representative with $|Z|\le M$ for
products with measures. Suppose, in distributions, that
\begin{equation}\label{n3:existence:natural-equation}
 \operatorname{div}(aDZ+b)=e,
\end{equation}
where $a$ is measurable and symmetric, $\lambda I\le a\le\Lambda I$
for $0<\lambda\le\Lambda$, $|b|\le B$, and $e$ is a signed measure
satisfying
\begin{equation}\label{n3:existence:natural-source}
 |e|\le C_0(1+|DZ|^2)\dd x+\mu,\qquad
 \mu(B_r(x))\le C_\mu r^{n-2+\beta},\qquad \beta>0.
\end{equation}
Here $\mu$ is a positive measure, and its bound holds at every center
and all sufficiently small radii in a larger patch. Assume the weak
exponential product rule \eqref{n3:existence:exponential-identity} below
is valid with the specified representative. This is true, in particular,
for a function continuous across a flat interface and smooth on its
two closed sides, when the equation includes the normal-flux jump.
Then the Sobolev class of $Z$ has a H\"older representative on smaller
balls, with constants depending only on the displayed bounds and patch
separation. If $Z$ is already continuous, this is a uniform estimate
for that function. The conclusion applies up to a smooth face whenever
its reflected equation satisfies these hypotheses.
\end{lemma}

\begin{proof}
The proof has two steps. An exponential change absorbs the quadratic
term in $DZ$; a small potential correction then permits weak Harnack
to decrease the essential oscillation. Both steps work in the two
dimensions stated in the lemma. Restricting all radii to at most one,
we may replace $\beta$ by $\min\{\beta,1\}$ in the measure bound.
We retain the symbol $\beta$ for this positive reduced exponent;
this also makes the small-scale mollification bound below immediate.

Write $\mathcal L=-\operatorname{div}(aD)$ and, for $s\in\{-1,1\}$,
set $V_s=e^{skZ}$. The assumed product rule is
\begin{equation}\label{n3:existence:exponential-identity}
 \mathcal L V_s
 =\operatorname{div}(skV_sb)-skV_se
       -k^2V_s\bigl(aDZ\cdot DZ+b\cdot DZ\bigr).
\end{equation}
The term $V_se$ uses the fixed bounded representative. For a continuous
piecewise smooth function the identity follows by integration on the
two sides; its common trace multiplies the flux jump. Choose
$k\ge\max\{1,4C_0/\lambda\}$ and use
$B|DZ|\le(\lambda/2)|DZ|^2+B^2/(2\lambda)$.
The gradient-square term from $e$ is absorbed by the negative quadratic
term for both signs. Since $e^{-kM}\le V_s\le e^{kM}$, this gives
\begin{equation}\label{n3:existence:exponential-subsolution}
 \mathcal L V_s\le\operatorname{div}F_s+\nu,
 \qquad |F_s|\le C,
 \qquad \nu=C(\dd x+\mu)\ge0.
\end{equation}
Thus no small-ball estimate for $|DZ|^2$ is needed.

On a ball $B_R$ in the smaller patch solve the zero-Dirichlet problem
$\mathcal L w_s=\operatorname{div}F_s+\nu$.
For the bounded vector source, the scalar bounded-solution estimate
with any fixed $p>n$ gives
\[
 \|w_{F_s}\|_\infty
 \le CR^{1-n/p}\|F_s\|_{L^p(B_R)}\le CR;
\]
this is the scaled divergence-source estimate of
\cite[Theorem~2.6, p.~50]{LittmanStampacchiaWeinberger1963}.
For the positive measure source, extend $a$ elliptically to $B_{2R}$.
The scalar Green comparison theorem
\cite[Theorem~7.1 and Remark~2, p.~66]{LittmanStampacchiaWeinberger1963},
and domain monotonicity, give
\[
 0\le\mathcal G_{B_R}(x,y)\le\mathcal G_{B_{2R}}(x,y)
                    \le C|x-y|^{2-n},\qquad x,y\in B_R.
\]
Using $B_{2R}$ keeps both points in a fixed compact interior region
of the comparison ball after rescaling. Dyadic annuli and
\eqref{n3:existence:natural-source} imply
\begin{equation}\label{n3:existence:green-morrey-correction}
 \begin{split}
 \sup_{x\in B_R}\int_{B_R}\mathcal G_{B_R}(x,y)\dd\nu(y)
 &\le C\sum_{j\ge0}(2^{-j}R)^{2-n}
                       \nu(B_{2^{1-j}R}(x))\\
 &\le C(R^2+R^\beta).
 \end{split}
\end{equation}
All balls used in this estimate lie in the chosen larger patch.

This potential also gives an energy solution. Mollify the restriction
of $\nu$ to $B_R$ and restrict the resulting positive density again to
$B_R$. Its ball-growth bounds are uniform: above the smoothing scale
use an enlarged ball, and below it use the mollified density bound.
The variational solutions have the preceding supremum bound and obey
\[
 \lambda\int|Dw_j|^2\le\int w_j\dd\nu_j
              \le\|w_j\|_\infty\nu_j(B_R).
\]
Weak compactness in $W^{1,2}_0$ and distributional convergence produce
the desired measure-source solution. Adding the vector-source solution,
and taking $R\le1$, yields
\begin{equation}\label{n3:existence:correction-size}
 \|w_s\|_\infty\le\varepsilon_R,
 \qquad\varepsilon_R=C(R+R^\beta).
\end{equation}

All extrema in the rest of this proof are essential extrema. Put
$M_R=\operatorname*{ess\,sup}_{B_R}Z$ and
$m_R=\operatorname*{ess\,inf}_{B_R}Z$. The functions
\begin{equation}\label{n3:existence:exponential-deficits}
 H_+=e^{kM_R}-e^{kZ}+w_++\varepsilon_R,
 \qquad
 H_-=e^{-km_R}-e^{-kZ}+w_-+\varepsilon_R
\end{equation}
are nonnegative almost everywhere and satisfy $\mathcal L H_\pm\ge0$.
At least half of $B_{R/2}$ satisfies either
$Z\le(M_R+m_R)/2$ or $Z\ge(M_R+m_R)/2$.
On the first set the positive exponential deficit is at least
$c(M_R-m_R)$; on the second set the negative exponential deficit
has that lower bound. Their derivatives are bounded above and away
from zero on $[-M,M]$.
Weak Harnack for the corresponding $H_\pm$ therefore either lowers
the essential supremum or raises the essential infimum on $B_{R/4}$
by at least $c'(M_R-m_R)-C\varepsilon_R$. Hence
\begin{equation}\label{n3:existence:oscillation-contraction}
 \operatorname*{ess\,osc}_{B_{R/4}}Z
 \le(1-c')\operatorname*{ess\,osc}_{B_R}Z+C(R+R^\beta).
\end{equation}
Choose a positive exponent smaller than $\min\{1,\beta\}$ and than
$-\log(1-c')/\log4$. Iteration gives an essential oscillation bound
$Cr^\gamma$ at every center in a smaller patch.
The averages of $Z$ on balls centered at $x$ are consequently Cauchy
as their radii decrease to zero. Their limit defines $Z^*(x)$ at
every center; comparison on overlapping balls gives
$|Z^*(x)-Z^*(y)|\le C|x-y|^\gamma$.
Lebesgue differentiation identifies $Z^*=Z$ almost everywhere.
This proves the representative assertion. For continuous $Z$ the
representatives agree everywhere. No conclusion about arbitrary
values of the initially chosen Borel representative on a null set
is required.

For the boundary application, flatten a face and set
$S=\operatorname{diag}(1,\ldots,1,-1)$. An even extension uses
$a_-=Sa_+S$ and $b_-=Sb_+$; its plane source is twice the
prescribed total normal flux. An odd extension after subtracting a
constant Dirichlet value uses $a_-=Sa_+S$, $b_-=-Sb_+$ and the
sign-reflected bulk source; its normal flux is continuous.
The coordinate volume density is included in these coefficients.
For the smooth reflected solutions used here the product rule holds,
and a bounded plane density has ball mass $O(r^{n-1})$, giving
$\beta=1$ for that part of the measure. This verifies the boundary
clause whenever the remaining reflected measure has the stated growth.
\end{proof}

\subsection{Boundary regularity for the coupled equations}

We return to the deformation. Its regularity must be obtained in
the order $Df$, then $Z$, then $D^2f$, and finally $DZ$. In the
last step, small oscillation of $Z$ absorbs its quadratic gradient
term. This order avoids assuming the Hessian estimate that the
second equation still needs.

\begin{proposition}[Classical bounds for bounded-variable solutions]
\label{n3:existence:classical-bounds}
Fix the deformation parameters and assume the bounds in
Proposition~\ref{n3:apriori:bounds}.  Every smooth solution in the
homotopy of Proposition~\ref{n3:deformation:homotopy}, with
$t\ge-\epsilon$, has uniform local classical estimates of every
finite order, up to the inner and outer faces.  At fixed $N$ the
constants are independent of the homotopy parameter and of all
sufficiently large outer radii $R$.  The local patches can be chosen
with a uniform positive radius throughout the end.
\end{proposition}

\begin{proof}
\emph{First, control $Df$ and the Hessian measure.}
Proposition~\ref{n3:apriori:bounds} bounds $f,Z,t$ and
$\sigma=|\nabla f|$. At fixed parameters this bounds $u,l,D$
above and away from zero and gives $\chi\ge\chi_0>0$.  After dividing the scalar trace
equation by $l/\sqrt D$, its right side is
\begin{equation}\label{n3:existence:normalized-f-source}
 G_f=\frac{\sqrt D}{l}
         \bigl(F_s-\tr_{A_\chi}K\bigr).
\end{equation}
It is bounded throughout the homotopy, including its added collar
terms.  Lemma~\ref{n3:existence:gradient} therefore gives a uniform
$C^{1,\alpha}$ estimate for $f$ and
\begin{equation}\label{n3:existence:f-morrey-application}
 \int_{B_r}|\Hess f|^2\le Cr^{1+2\alpha}
\end{equation}
on the uniform patches, including boundary patches where $f$ is
constant.

\emph{Second, control the oscillation of $Z$.}
After multiplication by the Riemannian volume density, the second
equation has coordinate principal coefficients
$a^{ij}=4\sqrt{\det g}\,uA_\chi^{ij}$, a bounded flux vector $b$, and a source
bounded in absolute value by
\begin{equation}\label{n3:existence:Z-quadratic-source}
 C\bigl(1+|DZ|^2+|D^2f|^2\bigr).
\end{equation}
Indeed $t$ is a smooth function of $x,Z,Df$ on the bounded variable
range.  Thus $Dt$ is a bounded linear combination of $DZ,D^2f$
and fixed smooth terms; the expression $\mathcal T$ is quadratic
in these quantities.  All other source terms have lower order.
The same volume density is included in $b,e$; it is smooth and bounded
above and away from zero on the chosen patches, so it preserves these
bounds and uniform ellipticity.  The prescribed conormal flux at the inner boundary is
bounded.  At the outer face $Z=0$.
Apply Lemma~\ref{n3:existence:natural-growth}, with
$\mu=C|D^2f|^2\dd x$ and the bounded reflected plane density.
Equation~\eqref{n3:existence:f-morrey-application} supplies its measure
growth hypothesis.  We obtain a uniform positive H\"older exponent
for $Z$.

\emph{Third, gain two derivatives of $f$.}
The coefficients and source of the scalar equation are smooth
functions of $x,Z,Df$ on the bounded range. They are therefore
H\"older now that $Z$ and $Df$ are H\"older.  Interior and constant-Dirichlet boundary Schauder
estimates give $f\in C^{2,\theta}$ for a uniform $\theta>0$.
Expand the $Z$-equation in nondivergence form.  Its principal
coefficients are uniformly elliptic and $C^{0,\theta}$, while its
right side is bounded by $C(1+|DZ|^2)$.
On an inner face $Df$ is normal.  Therefore $A_\chi$ has zero
normal--tangential components there and normal eigenvalue $\chi$.
Its boundary equation solves for
\begin{equation}\label{n3:existence:ordinary-normal-condition}
 \partial_\nu Z=\mathcal H_s(x,Z,f,Df),
\end{equation}
where $\mathcal H_s$ is smooth on a fixed bounded variable range.
Its derivatives with respect to $x$ and its displayed scalar
arguments are bounded at fixed parameters.  The outer condition
remains constant Dirichlet.  These are different, disjoint smooth
boundary components; there is no change of boundary condition along
an edge or a corner.

\emph{Fourth, absorb the quadratic growth and control $DZ$.}
Continuity of $Z$ supplies the small factor needed for this step.
Fix $p>3$.  The linear local $W^{2,p}$ estimate for the leading
nondivergence operator with Equation~\eqref{n3:existence:ordinary-normal-condition}
or outer Dirichlet data has the form, on fixed nested patches
$U_x\Subset V_x$,
\begin{equation}\label{n3:existence:Wp-pre-absorption}
 \|Z\|_{W^{2,p}(U_x)}
 \le C\left(1+\|DZ\|_{L^{2p}(V_x)}^2
                 +\|Z\|_{W^{1,p}(V_x)}\right).
\end{equation}
Here the boundary norm is controlled by
\begin{equation}\label{n3:existence:boundary-trace-composition}
 \|\mathcal H_s(x,Z,f,Df)\|_{W^{1-1/p,p}(\partial V_x)}
 \le C\bigl(1+\|Z\|_{W^{1,p}(V_x^+)}\bigr),
\end{equation}
using trace and composition after extending its smooth expression
through the collar.  The derivatives of $f,Df$ entering this
extension are already bounded.  The leading matrix has a fixed
H\"older modulus, so its small oscillation on sufficiently small
patches permits the usual freezing and absorption in the linear
estimate.  The normal boundary operator is uniformly oblique and
satisfies the complementing condition. The general boundary theory is
developed in the two papers of Agmon, Douglis and Nirenberg
\cite{AgmonDouglisNirenberg1959,AgmonDouglisNirenberg1964}. The scalar
estimate used here is
\cite[Theorem~15.2, estimate~(15.5), p.~704]{AgmonDouglisNirenberg1959}.

The outer patches $U_x,V_x,V_x^+$ and the constant in
\eqref{n3:existence:Wp-pre-absorption} are fixed before the next radius
$h$ is chosen; only the interpolation is localized further.
Let $[Z]_{C^{0,\theta}}\le H_0$. On balls or half-balls of
radius $h$, the scaled Gagliardo--Nirenberg inequality, applied after
subtracting a constant, gives
\begin{equation}\label{n3:existence:small-oscillation-GN}
 \begin{split}
 \|DZ\|_{L^{2p}(B_h)}^2
 &\le C\operatorname{osc}_{B_h}Z\,
                    \|D^2Z\|_{L^p(B_h)}\\
 &\quad+Ch^{3/p-2}(\operatorname{osc}_{B_h}Z)^2.
 \end{split}
\end{equation}
One can use a bounded extension operator on a half-ball; its constants
are uniform in the boundary charts.  Here the interpolation exponents are spatial dimension $3$,
derivative orders $2$ and $1$, top-derivative exponent $p>3$,
height exponent $\infty$, and interpolation parameter $1/2$;
the resulting gradient exponent is $2p$. This is
\cite[Lecture~II, Equation~(2.2), p.~125, and Remark~5, p.~126]{Nirenberg1959},
including the lower-order term on a bounded patch. Subtracting a
constant gives the oscillation, rather than the absolute height,
as the coefficient of $\|D^2Z\|_p$.
Cover $V_x$ by radius-$h$ patches of bounded overlap and take the
$\ell^p$ sum of Equation~\eqref{n3:existence:small-oscillation-GN}.
It follows that
\begin{equation}\label{n3:existence:global-local-GN}
 \|DZ\|_{L^{2p}(V_x)}^2
 \le CH_0h^\theta\|D^2Z\|_{L^p(V_x^+)}
           +CH_0^2h^{2\theta-2}.
\end{equation}
Also, for any $\eta>0$,
$\|Z\|_{W^{1,p}(V_x)}
 \le\eta\|D^2Z\|_{L^p(V_x^+)}+C_\eta$.
Each enlarged patch $V_x^+$ is covered by a bounded number of the
patches $U_y$, independently of $R$.  Set
$S=\sup_y\|Z\|_{W^{2,p}(U_y)}$, finite for each smooth finite-domain
solution.  Equations~\eqref{n3:existence:Wp-pre-absorption} and
\eqref{n3:existence:global-local-GN} give
\begin{equation}\label{n3:existence:uniform-local-absorption}
 S\le C(H_0h^\theta+\eta)S+C(h,\eta).
\end{equation}
Choose $h$ and then $\eta$ so the first coefficient is less than
$1/2$.  We obtain a uniform $W^{2,p}$ bound and hence a
$C^{1,1-3/p}$ bound for $Z$.

\emph{Finally, obtain all remaining derivatives.}
The bounds $f\in C^{2,\theta}$ and $Z\in C^{1,1-3/p}$ make the
expanded $Z$ source H\"older and its normal datum
$\mathcal H_s(x,Z,f,Df)$ of class $C^{1,\theta'}$, after reducing
the positive exponent. Its principal coefficients are H\"older as
well, so both the interior and normal-boundary Schauder estimates apply.
First the nondivergence $Z$-equation has H\"older source, giving
$Z\in C^{2,\theta'}$; the scalar $f$-equation then gains another
derivative.  Repetition gives every finite classical derivative
supported by the smooth data.  The same procedure holds on the
uniform unit patches of the asymptotic end and on the large outer
spheres, whose normal-coordinate constants are uniform.  This
proves the asserted independence of $R$.  The prepared background has end symbol bounds of every order by
Definition~\ref{n3:prepared:data}; no regularity constant here
requires additional derivatives of the original weak-decay data.
\end{proof}

\subsection{Solving the scalar equation for arbitrary inputs}

The previous proposition estimates solutions of the coupled system.
To construct such a solution by degree, we first need an operator
defined on all of a Banach space of trial functions $Z$. In
particular, $t\ge-\epsilon$ cannot be assumed at this stage.
The scalar equation remains solvable because its loss of axial
ellipticity at large slope is at most reciprocal-linear.

\begin{lemma}[Scalar solve for unrestricted bounded inputs]
\label{n3:existence:scalar-solve}
Fix $N\ge4$, $\tau>0$, a finite smooth truncation $\Omega_R$,
and a homotopy parameter $s\in[0,4]$.  For every
$Z\in C^{1,b}(\overline{\Omega_R})$, $0<b<1$, the scalar trace
equation of Proposition~\ref{n3:deformation:homotopy}, with its
prescribed componentwise constant Dirichlet values for $f$, has a
unique solution.  It lies in $C^{3,b}$ and is smooth to the extent
allowed by the input.  The solution operator
\begin{equation}\label{n3:existence:scalar-operator}
 f=\mathcal S_s(Z)
\end{equation}
is continuous in $(s,Z)$, including across the concatenation points
of the homotopy.  On bounded subsets of $C^{1,b}$ its $C^{3,b}$
norm is uniformly bounded at these fixed parameters.  No lower
bound for $t$ is required.
\end{lemma}

\begin{proof}
For the scalar continuation below, use the affine space of
$C^{3,b}$ functions with the assigned Dirichlet values and the equation
space $C^{1,b}$. Its individual solutions meet the $C^3$ hypothesis of
Lemma~\ref{n3:existence:gradient}; no uniform third-derivative bound is
assumed before that lemma is applied.
The monotonicity $\partial Z/\partial t=1+pv/4>0$ defines
$t=t(x,Z,Df)$ smoothly for every finite input, including $Df=0$.
The principal matrix has no dependence on the value of $f$.
The right side is increasing in $h=\tau f$, with derivative
at least one. At a height extremum the collar terms vanish and
$Df=0$, so the equation gives fixed constant upper and lower
barriers. Including all prescribed boundary heights, we obtain
\begin{equation}\label{n3:existence:scalar-height}
 |h|\le C_0,
\end{equation}
where the constant can be chosen independently of $Z$ and the
homotopy parameter.  We include the fixed boundary values in its
choice.

\emph{Large slopes.}
Only $|Z|\le M$ is assumed here. As $\sigma\to\infty$,
the implicit equation forces
$t\to-\infty$ uniformly in $Z$.  There $p=N$ and $l=e^{Nt}$,
so its exact form is
\begin{equation}\label{n3:existence:implicit-large-slope}
 e^{8Z}=e^{8t}+e^{(2N+8)t}\sigma^2.
\end{equation}
Consequently
\begin{equation}\label{n3:existence:large-slope-asymptotics}
 \begin{split}
 t&=-\frac{\log\sigma}{N+4}+O(1),\qquad
 l\asymp\sigma^{-N/(N+4)},\\
 d&\asymp\sigma^{-8/(N+4)},\qquad
 \chi\asymp\sigma^{-8/(N+4)},\qquad
 \frac{\sqrt D}{l}\asymp\sigma.
 \end{split}
\end{equation}
The constants are uniform for $|Z|\le M$ at fixed $N$.
Since $N\ge4$, the exponent $8/(N+4)$ is at most one.  Combining
the large-slope estimates with compactness of the remaining variable
range yields
\begin{equation}\label{n3:existence:scalar-growth}
 \chi\ge\frac{c}{1+\sigma},\qquad
 |G_f|\le C(1+\sigma)
\end{equation}
whenever Equation~\eqref{n3:existence:scalar-height} holds.

\emph{Boundary slopes.}
The estimate $\chi\ge c/(1+\sigma)$ suggests a concave logarithmic
barrier: its axial second derivative can dominate terms of linear
growth in its slope.
Let $\delta$ be smaller than a fixed collar radius, half the distance
between distinct boundary components, and a normal injectivity
radius of a smooth extension of the finite domain.  On a boundary
distance collar use
\begin{equation}\label{n3:existence:logarithmic-modulus}
 \psi(r)=k^{-1}\log(1+Br),\qquad
 \psi''=-k(\psi')^2.
\end{equation}
The upper and lower barriers are the assigned face value plus or
minus $\psi$.  On these tests the axis is normal to the distance
leaves, so their principal trace is
$\chi\psi''+\psi'\Delta r$.  At slopes at least one,
Equation~\eqref{n3:existence:scalar-growth} gives
\begin{equation}\label{n3:existence:logarithmic-dominance}
 -k\chi(\psi')^2\le-\tfrac12ck\psi'.
\end{equation}
Choose $k$ large enough to dominate all the bounded geometry terms
and the linear growth in Equation~\eqref{n3:existence:scalar-growth}.
Next shrink $\delta$ so that $1/(2k\delta)$ exceeds any fixed
large-slope threshold used here.  Finally choose $B$ large enough
that $B\delta\ge1$ and $\psi(\delta)>4C_0/\tau$.
The latter exceeds twice the bound for $\operatorname{osc}f$
from \eqref{n3:existence:scalar-height}.  Then $\psi'\ge1/(2k\delta)$ on this collar and
the two barriers compare by the height monotonicity.  This gives
a uniform boundary gradient bound.

\emph{Interior slopes.}
Use the same modulus for differences at two points. After increasing
$B$ if needed, consider
\begin{equation}\label{n3:existence:doubling-function}
 f(x)-f(y)-\psi(\operatorname{dist}(x,y)),
 \qquad 0<\operatorname{dist}(x,y)<\delta.
\end{equation}
The distance is taken in the fixed smooth extension.  The boundary barriers exclude a positive maximum with one endpoint
on a face: the other point lies in that face's collar, its distance
to the face is at most the distance between the pair, and $\psi$ is
increasing. The separation choice prevents endpoints on distinct faces.  The inequality $\psi(\delta)>2\operatorname{osc}f$
excludes the other boundary of the two-point domain.  At a positive
interior maximum let $r=\operatorname{dist}(x,y)$ and
$q=\psi'(r)$.  The two gradients have length $q$ and are parallel
transports along the connecting geodesic.  Vary both endpoints
simultaneously in parallel transverse directions.  The second
variation formula for the short geodesic yields
\begin{equation}\label{n3:existence:two-point-transverse}
 \tr_{e_x^\perp}\Hess f(x)
 -\tr_{e_y^\perp}\Hess f(y)\le Cqr.
\end{equation}
Varying each endpoint separately along that geodesic gives
$f_{;e_xe_x}(x)\le\psi''(r)$ and
$f_{;e_ye_y}(y)\ge-\psi''(r)$.  Subtracting the two scalar
equations therefore gives
\begin{equation}\label{n3:existence:two-point-contradiction}
 -2C(1+q)
 \le G_f(x)-G_f(y)
 \le Cqr+(\chi_x+\chi_y)\psi''(r).
\end{equation}
By Equations~\eqref{n3:existence:scalar-growth} and
\eqref{n3:existence:logarithmic-modulus}, the last term is at most
$-ckq$ at the large comparison slopes.  Our choice of $k$ gives
a contradiction.  Thus Equation~\eqref{n3:existence:doubling-function}
is nonpositive.  Reversing $x,y$ proves the same bound for the
absolute difference and, on letting $y\to x$, a global Lipschitz
bound for $f$.  The endpoint variations are interior variations in the smooth
extension; the connecting short geodesic need not stay inside
$\Omega_R$. This argument used no continuity modulus for
$\chi_x-\chi_y$; only the common transverse eigenvalues and the
radial lower bound were needed.

\emph{Continuation and uniqueness.}
The height and gradient estimates supply compactness for a scalar
continuation. Interpolate the normalized equation with
$\Delta f=\tau f$, retaining its prescribed boundary values.  Explicitly,
for $a\in[0,1]$ take principal matrix
$(1-a)I+aA_\chi$ and right side
$(1-a)\tau f+aG_f$.  The height sign, the linear gradient growth,
and Equation~\eqref{n3:existence:scalar-growth} persist uniformly in
$a$.  The height and logarithmic gradient bounds just proved are
therefore uniform.  After these bounds the interpolating equations
are uniformly elliptic.  Their matrices still have the axial form,
with axial eigenvalue $(1-a)+a\chi$, so
Lemma~\ref{n3:existence:gradient} applies.  Schauder estimates then
give a uniform $C^{2,\theta}$ bound at fixed input $Z$.
The coefficients and right side are now Lipschitz in position on
bounded input sets, which permits their $C^{0,b}$ estimates and
an upgrade to $C^{2,b}$.  Differentiating once, or using the next
Schauder estimate, yields $f\in C^{3,b}$ because $Z\in C^{1,b}$.

The scalar linearization has the form
$a^{ij}D_{ij}\varphi+b^iD_i\varphi-c\varphi$ with
$c\ge c_1\tau>0$ on the bounded variable range.  Indeed the
principal matrix is independent of the height, while
$\partial_fG_f\ge\tau\sqrt D/l>0$.
The Dirichlet maximum principle gives uniqueness and a zero kernel;
the linear elliptic Dirichlet theory gives an inverse from
$C^{1,b}$ to $C^{3,b}_0$, where the subscript denotes zero data on
the entire boundary.  The
implicit-function theorem gives openness of the interpolation
parameter set.  The uniform classical bounds and compactness give
closedness, starting from the invertible Laplace endpoint.  This
proves existence.  The same maximum principle applied to the
difference of two solutions proves uniqueness of the nonlinear
solution.  The implicit-function theorem, or compactness and
uniqueness, gives continuous dependence on $Z$ and on the
homotopy parameter.  At the finitely many concatenation points the
scalar equations agree, so the same continuity holds there.
\end{proof}


\subsection{Degree on the admissible set and exhaustion}

We can now construct the compact map. The homotopy parameter
$s\in[0,4]$ has a separate role from the deformation integer $N$:
in Proposition~\ref{n3:deformation:homotopy}, $s=0$ is the original
system, $s=1$ removes the drift, $s=2$ completes the collar
modification, and $s=3$ makes the scalar solution $f$ identically
zero. On $[3,4]$ the divergence source and the inner flux are
multiplied together by $4-s$. The final compact map is consequently
the zero map. The remaining task is to show that no fixed point
can cross the boundary of the admissible set along this path.

\begin{lemma}[Degree on varying open sets]\label{b:lem:moving-degree}
Let $I=[a,b]$, let $X$ be a real Banach space, and let
$T:I\times X\to X$ be continuous and map bounded sets to relatively
compact sets. Let $\mathcal U\subset I\times X$ be bounded and
relatively open. Suppose no fixed point $Z=T(s,Z)$ belongs to the
relative boundary of $\mathcal U$. Then, for
$\mathcal U_s=\{Z:(s,Z)\in\mathcal U\}$, the Leray--Schauder degree
\[
 \deg(I_X-T_s,\mathcal U_s,0)
\]
is well defined and independent of $s$. Empty sections have degree zero.
If $T_b$ is the zero map and $0\in\mathcal U_b$, the common degree
is one, and every section contains a fixed point.
\end{lemma}

\begin{proof}
The fixed points in $\overline{\mathcal U}$ form a compact set
$\mathfrak F$. Indeed, boundedness of $\mathcal U$, compactness of $I$
and of the images of $T$ give a convergent subsequence of any sequence
of such fixed points; continuity preserves the fixed-point equation.
The hypothesis places $\mathfrak F$ inside $\mathcal U$.
Also $\{s\}\times\partial\mathcal U_s$ lies in the relative boundary
of $\mathcal U$, so each section degree is defined.

Fix $s_0\in I$. Cover its compact fixed-point slice by finitely many
open balls in $X$ whose closed balls remain in $\mathcal U_s$ for
all $s$ sufficiently close to $s_0$. Relative openness and the finite
cover permit this common parameter interval. Let $U$ be their union.
All fixed points in the closed sections for these nearby parameters
lie in $U$ after shortening the interval: otherwise compactness of
$\mathfrak F$ would give an omitted fixed point in the slice at $s_0$.
Excision identifies the degree on each $\mathcal U_s$ with that on
$U$, and fixed-domain homotopy invariance makes the latter constant.
If the slice is empty, compactness instead gives no nearby fixed
points in the closed sections, so those degrees are all zero.
Thus the section degree is locally constant, and connectedness of $I$
makes it constant. These are the excision and homotopy properties of
Leray--Schauder degree
\cite[Sections~12--14, pp.~59--61]{LeraySchauder1934}; see also \cite{Deimling1985}.
At the stated terminal map, normalization gives degree one; the
existence property of nonzero degree proves the last assertion.
\end{proof}

\begin{theorem}[Existence on finite truncations]
\label{n3:existence:deformation}
For the prepared domain of Proposition~\ref{n3:domains:prepared-domain}
and every sufficiently small fixed $\epsilon>0$, there is an
integer $N_0\ge4$ and, for each $N\ge N_0$, a radius
$R_{\min}(N)$, with $R_{\min}(N)\to\infty$, such that the original
deformation system has a smooth solution on every $\Omega_R$ with
$R\ge R_{\min}(N)$.  It satisfies $t>-\epsilon$ on the closure.
At fixed $N$, these solutions have the local bounds of
Proposition~\ref{n3:existence:classical-bounds}, uniformly as
$R\to\infty$.  In particular, any sequence $R_j\to\infty$ of such radii has a
subsequence converging smoothly on compact subsets of
$\overline\Omega$ to a solution of the original equations with
$t\ge-\epsilon$.
\end{theorem}

\begin{proof}
\emph{Choose the parameters before defining degree.}
Fix the prepared data, thresholds, collars and $\epsilon$.
Choose an integer $N_0$ so that
\begin{equation}\label{n3:existence:n-choice}
 N_0\ge\max\{4,\lfloor2/\epsilon\rfloor+1\}
\end{equation}
and so that the floor test and all estimates of
Proposition~\ref{n3:apriori:bounds} apply for $N\ge N_0$.
Let $R_{\mathrm{ap}}(N)$ dominate every lower-radius requirement
in the a priori estimates and in the floor exclusion; let
$R_{\mathrm{geo}}$ contain all compact boundary faces, collars and
supports of the construction.  The end barrier of
that proposition gives a constant $C_{\mathrm{out}}$, independent
of $N,R,s$, such that every admissible or floor-contact homotopy
solution has
\begin{equation}\label{n3:existence:outer-gradient}
 |\nabla f|\le
       C_{\mathrm{out}}\tau^{-1}R^{-1-\gamma}
       \quad\text{on the outer sphere}.
\end{equation}
The scalar last stage has $f=0$ and satisfies the same bound.
Choose, enlarging to a radius of a smooth truncation if necessary,
\begin{equation}\label{n3:existence:R-choice}
 \begin{split}
 R_{\min}(N)=\max\biggl\{&N,R_{\mathrm{geo}},R_{\mathrm{ap}}(N),\\
 &\left(\frac{2C_{\mathrm{out}}\ell}
               {\tau\sqrt{e^{8\epsilon}-1}}
        \right)^{1/(1+\gamma)}\biggr\}.
 \end{split}
\end{equation}
At $Z=0$, the value of its implicit expression at $t=-\epsilon$
is $-\epsilon+\frac18\log(1+\ell^2\sigma^2)$.
Equations~\eqref{n3:existence:outer-gradient} and
\eqref{n3:existence:R-choice} make this strictly negative.  Strict
monotonicity of that expression in $t$ therefore excludes an outer
floor contact.  The term $N$ in Equation~\eqref{n3:existence:R-choice}
ensures $R_{\min}(N)\to\infty$.

The sublevel and collar estimates apply to arbitrary smooth
admissible or floor-contact solution sequences with $N\to\infty$
and expanding outer cutoffs, uniformly over the compact ranges of
the homotopy parameters.  If the needed conclusion failed for
arbitrarily large $N$ with $R\ge R_{\min}(N)$, one could select
such a violating sequence.  The chosen radius convention would
make it precisely an expanding sequence excluded by those proofs.
Thus a single $N_0$, followed by the radius choices above, suffices
for every homotopy parameter.  No upper bound for $Z$ is inserted
in the earlier contact estimates: there $\sigma\asymp\tau^{-1}$
and $l\ge\ell$ alone give $d\lesssim(\tau/\ell)^2$.

\emph{Define the compact map on all bounded input sets.}
For the rest of the degree argument fix $N\ge N_0$ and
$R\ge R_{\min}(N)$. Work in the Banach space
\begin{equation}\label{n3:existence:banach-space}
 X_b=\{Z\in C^{1,b}(\overline{\Omega_R}):
                       Z=0\text{ on the outer face}\},
 \qquad 0<b<1.
\end{equation}
For an input $Z\in X_b$, first solve
$f=\mathcal S_s(Z)$ by Lemma~\ref{n3:existence:scalar-solve}.
Compute $t,u,A_\chi,\mathcal T$ and the full homotopy source
from this input pair.  Write $B_s$ for the frozen drift summand in
the flux, $E_s$ for the full right side of its divergence equation,
and $q_s$ for the prescribed inner flux.  Thus before the final
scaling $B_s=\lambda_suA_\chi K(w,\cdot)$ and $E_s=u\Xi$;
on the final interval $E_s$ and $q_s$ have the simultaneous scalar
factor prescribed in Proposition~\ref{n3:deformation:homotopy}.
Define $\mathcal K_s(Z)=\widetilde Z$ by the linear mixed problem
\begin{equation}\label{n3:existence:compact-map-equation}
 \begin{cases}
  \Div(4uA_\chi\nabla\widetilde Z+B_s)=E_s
                           &\text{in }\Omega_R,\\
  (4uA_\chi\nabla\widetilde Z+B_s)_\nu=q_s
                           &\text{on }B,\\
  \widetilde Z=0           &\text{on the outer face}.
 \end{cases}
\end{equation}
All coefficients in this problem are computed from the input
$(\mathcal S_s(Z),Z)$; the only unknown is $\widetilde Z$.
The normal $\nu$ on the inner faces points into the domain.
For a test function $\varphi$ vanishing on $S_R$, the weak form is
\[
 \int_{\Omega_R}4uA_\chi\nabla\widetilde Z\cdot\nabla\varphi
 =-\int_{\Omega_R}E_s\varphi
  -\int_{\Omega_R}B_s\cdot\nabla\varphi
  -\int_B q_s\varphi.
\]
The sign of the boundary term comes from using $-\nu$ as the
integration normal. The scalar estimate bounds the leading
coefficient above and away from zero on each bounded input set.
The nonempty outer Dirichlet face supplies Poincar\'e's inequality,
so the bilinear form is coercive and has a unique solution.

On bounded input sets in $C^{1,b}$,
Lemma~\ref{n3:existence:scalar-solve} bounds $f$ in $C^{3,b}$.
Hence the leading coefficient and drift in
Equation~\eqref{n3:existence:compact-map-equation} are $C^{1,b}$,
the bulk source is $C^{0,b}$, and the prescribed boundary data are
$C^{1,b}$.  The derivatives of $t(x,Z,Df)$ involve only $DZ$
and $D^2f$, so this assertion includes the quadratic expression
$\mathcal T$.  Linear regularity gives a bounded output set in
$C^{2,b}$.  The inclusion $C^{2,b}\hookrightarrow C^{1,b}$ is
compact on the finite domain.  Continuous dependence of the scalar
solve and the linear problem makes $(s,Z)\mapsto\mathcal K_s(Z)$
a continuous compact homotopy on bounded sets.  These assertions
hold also at the concatenation parameters because the defining
equations coincide there.  A fixed point solves the coupled system;
repeated Schauder estimates make it smooth.

\emph{Exclude fixed points from the boundary of admissibility.}
Write $t_s[Z]=t(x,Z,D\mathcal S_s(Z))$. The scalar solution must
be computed before this admissibility test. Define
\begin{equation}\label{n3:existence:admissible-open-set}
 \mathcal O=\{(s,Z)\in[0,4]\times X_b:
       \min_{\overline{\Omega_R}}t_s[Z]>-\epsilon,
       \ \|Z\|_{C^{1,b}}<M_1\}.
\end{equation}
It is relatively open because the scalar solve and the implicit
function defining $t$ depend continuously on $(s,Z)$.
Propositions~\ref{n3:apriori:bounds} and
\ref{n3:existence:classical-bounds} give a uniform bound in $C^{1,b}$
for all fixed points satisfying $t\ge-\epsilon$.  If initially
needed, use their higher classical bounds to reach the chosen
exponent $b$.  Choose $M_1$ larger than this bound.  No fixed point
lies on the norm boundary of $\mathcal O$.  No fixed point lies on
its floor boundary either: outer contacts were excluded above and
all remaining contacts are excluded by
Proposition~\ref{n3:deformation:floor}.

\emph{Apply degree on the varying admissible sections.}
The sections $\mathcal O_s$ need not be constant in $s$.
Lemma~\ref{b:lem:moving-degree} applies with $I=[0,4]$, $X=X_b$,
$T_s=\mathcal K_s$ and $\mathcal U=\mathcal O$.
We have proved continuity and compactness on bounded input sets;
the norm condition makes $\mathcal O$ bounded and the scalar solve
makes it relatively open. Its relative boundary can contain only
norm contact or floor contact, and neither can be a fixed point by
the preceding estimates. Thus its section degree is constant.

At $s=4$, the scalar equation has the unique solution $f=0$.
The drift vanishes, and both $E_4$ and $q_4$ vanish for every input.
Equation~\eqref{n3:existence:compact-map-equation} then has output
identically zero.  Moreover, $t_4[0]=0>-\epsilon$, so
$0\in\mathcal O_4$ and
\begin{equation}\label{n3:existence:degree-one}
 \deg(I-\mathcal K_4,\mathcal O_4,0)=1.
\end{equation}
Homotopy invariance gives the same degree at $s=0$.  There is
therefore a fixed point in $\mathcal O_0$, which is the desired
smooth solution with $t>-\epsilon$.

\emph{Remove only the outer truncation.}
Keep $N$ and all other deformation parameters fixed.
Proposition~\ref{n3:existence:classical-bounds} gives estimates
independent of $R$ on an exhaustion by compact sets, using boundary
charts at the fixed inner faces. A diagonal subsequence converges
smoothly on compact subsets of $\overline\Omega$.
Both equations and their inner boundary conditions pass to the
limit. The strict finite-domain inequality gives the non-strict
limit inequality $t\ge-\epsilon$. The following section establishes
end normalization, decay and mass flux; those conclusions do not
follow from compact exhaustion alone.
\end{proof}
